% !TeX root = ../../Book4.tex
% 纲要标题：## 第7章　链同伦、映射锥与拟同构
% 纲要标签：b4:ch:06
% 纲要位置：计划纲要.md，第 3342 行。

\chapter{链同伦、映射锥与拟同构}
\label{b4:ch:06}
\BookChapterNumber{b4:ch:06}

\begin{chapterabstract}
链同伦保证两个复形映射在同调上诱导相同映射；同伦等价则比拟同构要求更多。本章通过映射锥和映射柱比较这些概念，再用有限消元证明双复形总化的比较结果，并说明无界总化时直和与直积的差别。
\end{chapterabstract}

\begin{learninggoals}
\begin{itemize}
\item 构造链同伦，判断正合复形何时可缩，并区别同调相同与映射同伦。
\item 写出映射锥、映射柱及其收缩同伦，检验移位和连接箭头的符号。
\item 用锥的正合性识别拟同构，说明拟同构何时能加强为同伦等价。
\item 正确总化双复形，使用有限消元证明正合性，并识别无界总化的不同结果。
\end{itemize}
\end{learninggoals}

把一条短正合列视为只在三个连续次数非零的复形，正合性说明它的同调处处为零。然而，要把恒等链映射写成「微分与一个降次映射的复合之和」，需要逐项找到能退回前一项的映射。例如，对交换群短正合列 \(0\rightarrow\mathbb Z\xrightarrow{2}\mathbb Z\rightarrow\mathbb Z/2\mathbb Z\rightarrow0\)，这样的降次映射会迫使商映射有截面，而这个截面并不存在。可缩正是复形与零复形同伦等价，因此一定零调；这个例子表明，零调并不保证可缩。同调计算无法区分该复形与零复形，二者却不同伦等价。对一般复形映射，也可以追问它在同调上成为同构以后，能否在复形层面找到同伦逆。映射锥把前一种可逆性转化为一个新复形的正合性，使两种要求可以直接比较。

\ProseParagraphBreak
本章使用\chapref{b4:ch:05}的复形运算及上同调长正合列。同伦与映射锥可参阅 Weibel\cite[Section 1.4, Lemma 1.4.5; Section 1.5, Lemma 1.5.3, Corollary 1.5.4]{Weibel1994HomologicalAlgebra}；双复形、总化与符号见\cite[Example 1.2.4, Sign Trick 1.2.5, Total Complexes 1.2.6]{Weibel1994HomologicalAlgebra}。本章使用的锥矩阵和移位按本书约定书写，与上述文献比较时需留意分量次序与符号的转换。

% !TeX root = ../../Book4.tex
% 纲要标题：### 7.1 链同伦
% 纲要标签：b4:sec:0601
% 纲要位置：计划纲要.md，第 3344 行。

\section{链同伦}
\label{b4:sec:0601}

同调看不见的复形映射未必都能用同伦消去。先研究可显式写成 \(\mathrm{d}h+h\mathrm{d}\) 的差别，便能区分零伦、可缩与仅仅正合。

% 纲要内容（仅作写作计划，尚非教材正文）：
%
% * 同伦定义与同调不变性
% * 零伦映射及可缩复形
% * 正合复形未必可缩
%

\subsection{同伦与同调不变性}
\label{b4:subsec:0601:01}

两个复形映射可能逐项不同，却在同调上完全相同。最容易控制的一类差别是能够写成“先作微分再作某个降次映射，加上先降次再作微分”。这类差别在余圈上自动变成余边。

\ProseParagraphBreak
\cref{b4:def:cochain-complex,b4:def:complex-map}在交换范畴中定义了复形与复形映射；在加性范畴中，同样可以取对象族 \(C^n\) 和满足 \(\mathrm{d}_C^{n+1}\mathrm{d}_C^n=0\) 的态射 \(\mathrm{d}_C^n:C^n\to C^{n+1}\)，并要求复形映射与微分交换。这些定义只用到零态射和复合。同伦还用到态射的加减法，因而也可在加性范畴中讨论；核、像及同调对象则仍在交换范畴中使用。

\begin{definition}\label{b4:def:chain-homotopy}
设 \(\mathcal A\) 是加性范畴，\(C,D\) 为 \(\mathcal A\) 中的上链复形，\(f,g:C\to D\) 是复形映射。若有态射族 \(h^n:C^n\to D^{n-1}\)（\(n\in\mathbb Z\)），使
\[
f^n-g^n=\mathrm{d}_D^{n-1}h^n+h^{n+1}\mathrm{d}_C^n
\quad(n\in\mathbb Z),
\]
则称 \(f\) 与 \(g\) \term[lian-tonglun]{链同伦}[chain homotopic]，记为 \(f\simeq g\)，并称 \(h\) 为一个链同伦。链指标中同一个公式写为 \(f_n-g_n=\partial_{D,n+1}h_n+h_{n-1}\partial_{C,n}\)，其中 \(h_n:C_n\to D_{n+1}\)。
\end{definition}
\symidx{\(f\simeq g\)：复形映射之间的链同伦关系}

上链指标下同伦降一次，链指标下升一次，两者都与微分方向相反。在次数 \(n\)，两个复合
\[
\mathrm{d}_D^{n-1}h^n:C^n\longrightarrow D^n,
\qquad h^{n+1}\mathrm{d}_C^n:C^n\longrightarrow D^n
\]
分别先降次再升次、先升次再降次，因而都能与 \(f^n-g^n\) 相加比较。这里保留“链同伦”这一通用名称，不因为指标升次而另换概念。

\ProseParagraphBreak
对模复形，\cref{b4:def:hom-complex}中的 Hom 复形恰好解释了这个公式的来源。态射族 \(h\) 是次数 \(-1\) 的元素，其微分为
\[
\mathrm{d}_{\operatorname{Hom}}h
=\mathrm{d}_Dh-(-1)^{-1}h\mathrm{d}_C
=\mathrm{d}_Dh+h\mathrm{d}_C.
\]
所以 \(f\simeq g\) 就是说零次余圈 \(f-g\) 是 Hom 复形中的余边。

\begin{proposition}\label{b4:prop:homotopy-cohomology}
设 \(\mathcal A\) 是加性范畴。在任意两个固定的 \(\mathcal A\) 中的上链复形之间，同伦是复形映射的等价关系，且与加法、左右复合相容。若 \(\mathcal A\) 还是交换范畴，\(f,g:C\to D\) 为同伦的复形映射，则对每个 \(n\in\mathbb Z\)，有 \(H^n(f)=H^n(g)\)。
\end{proposition}
\begin{proof}
零同伦、\(-h\) 及同伦之和分别给出自反、对称和传递性。若 \(u:C'\to C\)、\(v:D\to D'\) 为复形映射，则
\[
v(f-g)u=\mathrm{d}_{D'}(vhu)+(vhu)\mathrm{d}_{C'};
\]
这证明复合相容性。两个差的同伦相加给出加法相容性。

再设 \(\mathcal A\) 是交换范畴，记 \(z_C^n,z_D^n\) 为余圈的核嵌入，\(j_D^n:B^n(D)\to Z^n(D)\) 为余边到余圈的嵌入，并写
\[
\mathrm{d}_D^{n-1}=z_D^nj_D^n\pi_D^{n-1},
\qquad \pi_D^{n-1}:D^{n-1}\twoheadrightarrow B^n(D).
\]
设 \(f_Z^n,g_Z^n:Z^n(C)\to Z^n(D)\) 是 \(f,g\) 在余圈上诱导的态射。把同伦公式右复合 \(z_C^n\)，含 \(h^{n+1}\mathrm{d}_C^n\) 的项消失，再消去单态射 \(z_D^n\)，得到
\[
f_Z^n-g_Z^n=j_D^n\pi_D^{n-1}h^nz_C^n.
\]
同调商映射 \(q_D^n:Z^n(D)\to H^n(D)\) 满足 \(q_D^nj_D^n=0\)，故 \(q_D^n(f_Z^n-g_Z^n)=0\)。由 \(q_C^n:Z^n(C)\to H^n(C)\) 为满态射，得到 \(H^n(f)-H^n(g)=0\)。在模中，这就是余圈 \(z\) 的两个像之差 \(f(z)-g(z)=\mathrm{d}_Dh(z)\) 是余边。
\end{proof}

因此对任意加性范畴 \(\mathcal A\)，都可以把同伦的映射视为同一个态射，得到\term[tonglun-fanchou]{同伦范畴}[homotopy category] \(K(\mathcal A)\)。对象仍是复形，态射则是复形映射的同伦类；上述相容性保证态射的复合与加法不依赖代表。\symidx{\(K(\mathcal A)\)：复形的同伦范畴}例如投射模构成的加性范畴也有同伦范畴，不必要求它本身是交换范畴。

\ProseParagraphBreak
若 \(C,D\) 是左 \(R\)-模复形，则
\[
\begin{aligned}
Z^0\operatorname{Hom}_R^\bullet(C,D)
&=\{C\to D\text{ 的复形映射}\},\\
B^0\operatorname{Hom}_R^\bullet(C,D)
&=\{\mathrm{d}_Dh+h\mathrm{d}_C\mid h\text{ 的次数为 }-1\}.
\end{aligned}
\]
右边第二个集合恰是零伦映射所成的子群。将零次余圈模掉这个子群，便是将两个同伦的复形映射视为相同，故
\[
\operatorname{Hom}_{K(R\text{-}\mathbf{Mod})}(C,D)
=H^0\operatorname{Hom}_R^\bullet(C,D),
\]
这把同伦类的定义与 Hom 复形的零次同调联系起来。

\subsection{零伦映射及可缩复形}
\label{b4:subsec:0601:02}

\begin{definition}\label{b4:def:contractible-complex}
设 \(\mathcal A\) 为加性范畴，\(C,D\) 为其中的上链复形，\(f:C\to D\) 为复形映射。若 \(f\) 与零映射同伦，则称 \(f\) 为\term[linglun-yingshe]{零伦映射}[null-homotopic map]。若 \(1_C\simeq0\)，则称 \(C\) 为\term[kesuo-fuxing]{可缩复形}[contractible complex]；满足 \(1_{C^n}=\mathrm{d}_C^{n-1}h^n+h^{n+1}\mathrm{d}_C^n\) 的态射族 \(h^n:C^n\to C^{n-1}\)（\(n\in\mathbb Z\)）称为收缩同伦。若存在复形映射 \(g:D\to C\)，使 \(gf\simeq1_C\)、\(fg\simeq1_D\)，则称 \(f\) 为\term[tonglun-dengjia]{同伦等价}[homotopy equivalence]。
\end{definition}

同伦等价所需的 \(g\) 是同伦逆，两个复合只要求与恒等映射同伦；复形同构的严格逆当然满足这一要求，反向却不必成立。可缩复形在交换范畴中必零调，因为恒等同态在同调上等于零。反向所缺的条件是各次数的短正合列能否分裂：分裂把 \(C^n\) 分成当前的余圈与下一次余边的选定提升，微分便只负责把后一个部分送到下一项。

\begin{proposition}\label{b4:prop:contractible-split-exact}
设 \(\mathcal A\) 是交换范畴，\(C\) 为其中的上链复形。它可缩，当且仅当它零调，且对每个 \(n\in\mathbb Z\)，短正合列
\[
0\longrightarrow Z^n(C)\xrightarrow{z_C^n}C^n
\xrightarrow{\widetilde{\mathrm{d}}_C^n} Z^{n+1}(C)\longrightarrow0
\]
均分裂；其中 \(z_C^n:Z^n(C)\to C^n\) 为核态射，\(\widetilde{\mathrm{d}}_C^n\) 由 \(\mathrm{d}_C^n=z_C^{n+1}\widetilde{\mathrm{d}}_C^n\) 定义，并在零调假设下为满态射。
\end{proposition}
\begin{proof}
若 \(1_C=\mathrm{d}h+h\mathrm{d}\)，则已知 \(C\) 零调。\(\widetilde{\mathrm{d}}_C^n\) 是微分的像分解中到 \(Z^{n+1}(C)=B^{n+1}(C)\) 的满态射；这里满的是这个分量，不是到整个 \(C^{n+1}\) 的微分。收缩同伦 \(h^{n+1}\) 将下一次余圈降回 \(C^n\)，因而截面的候选是
\[
s^{n+1}\coloneqq h^{n+1}z_C^{n+1}:Z^{n+1}(C)\longrightarrow C^n.
\]
把次数 \(n+1\) 的同伦公式右复合 \(z_C^{n+1}\)，由 \(\mathrm{d}_C^{n+1}z_C^{n+1}=0\) 得
\[
z_C^{n+1}=\mathrm{d}_C^nh^{n+1}z_C^{n+1}
=z_C^{n+1}\widetilde{\mathrm{d}}_C^nh^{n+1}z_C^{n+1}.
\]
消去单态射 \(z_C^{n+1}\)，得到
\[
\widetilde{\mathrm{d}}_C^n\bigl(h^{n+1}z_C^{n+1}\bigr)
=1_{Z^{n+1}(C)}.
\]
所以 \(s^{n+1}\) 确是 \(\widetilde{\mathrm{d}}_C^n\) 的截面。

反过来，选各短列的截面 \(s^{n+1}:Z^{n+1}(C)\to C^n\)。分裂短正合列给出的同构具体为
\[
\alpha^n=(z_C^n,s^{n+1}):Z^n(C)\oplus Z^{n+1}(C)\xrightarrow{\sim}C^n.
\]
第一分量由微分消去，第二分量经微分成为下一项的余圈嵌入，故在这些双积坐标中，微分为 \(\mathrm{d}(z,w)=(w,0)\)。定义 \(h^n(z,w)=(0,z)\in Z^{n-1}(C)\oplus Z^n(C)\)，则 \(\mathrm{d}h(z,w)=(z,0)\)、\(h\mathrm{d}(z,w)=(0,w)\)，相加等于恒等。这里的坐标公式表示双积的分量态射，因而计算在一般交换范畴中也成立。
\end{proof}

例如在次数 \(n,n+1\) 放同一对象 \(M\)，微分为 \(1_M\)，其余各项为零。只取 \(h^{n+1}=1_M\)，其余同伦分量为零，就得到收缩同伦。它的两个非零项可以很大，同调却为零；可缩所说的“消去”是通过显式同伦实现的，并不是逐项零对象。

\subsection{正合复形的非可缩例子}
\label{b4:subsec:0601:03}

\begin{example}\label{b4:exm:acyclic-not-contractible}
把交换群短正合列 \(0\to\mathbb Z\xrightarrow{2}\mathbb Z\to\mathbb Z/2\mathbb Z\to0\) 放在次数 \(0,1,2\)，便得到零调复形。若它可缩，记次数一的商映射为 \(\mathrm{d}^1\)，由于 \(\mathrm{d}^2=0\)，次数二的同伦公式就是 \(1_{\mathbb Z/2\mathbb Z}=\mathrm{d}^1h^2\)。这要求 \(h^2:\mathbb Z/2\mathbb Z\to\mathbb Z\) 为商映射的截面。但任何这样的群同态都为零，因此不存在截面，复形不可缩。
\end{example}

\begin{example}\label{b4:exm:acyclic-free-periodic}
令 \(R=\mathbb Z/4\mathbb Z\)，在每个整数次数取自由模 \(R\)，微分均为乘 \(2\)。因为核与像同为 \(2R\)，复形正合。每个 \(R\)-线性 \(h^n:R\to R\) 都是乘某个 \(a_n\in R\)。若有收缩同伦，便须
\[
1=2a_n+2a_{n+1}\quad\text{在}\;R\;\text{中成立},
\]
右端是偶数剩余类，矛盾。这说明即使各项自由，无界正合复形也未必可缩。

\ProseParagraphBreak
此时余圈 \(2R\cong R/(2)\) 不是投射模。否则满射 \(R\xrightarrow{2}2R\) 会有截面 \(s:2R\to R\)；由 \(2s(2)=s(2\cdot2)=0\) 可知 \(s(2)\in2R\)，所以乘 \(2\) 又将它送到零，无法得到所要求的 \(2\)。因此各项自由并没有保证上述余圈短正合列分裂。
\end{example}

作为对照，对任意含幺环 \(R\)，各项均为投射幺性左 \(R\)-模的上有界零调复形必可缩。取 \(b\) 使 \(n>b\) 时 \(C^n=0\)，则 \(Z^b(C)=C^b\) 投射；于是 \(0\to Z^{b-1}(C)\to C^{b-1}\to Z^b(C)\to0\) 分裂，且 \(Z^{b-1}(C)\) 作为投射模的直和因子仍投射。如此向较低次数归纳，所有余圈短正合列都分裂，再用上述命题得到可缩性。周期反例没有可供这一步归纳开始的最高次数。后面比较投射消解时，也将从一个有限端开始逐次提升同伦，不能在没有起点的无界复形上直接照搬。

% !TeX root = ../../Book4.tex
% 纲要标题：### 7.2 映射锥
% 纲要标签：b4:sec:0602
% 纲要位置：计划纲要.md，第 3350 行。

\section{映射锥}
\label{b4:sec:0602}

一个复形映射也可以成为新复形的微分组成部分。映射锥将它在上同调上诱导的映射的核与余核信息集中到长正合列中，柱构造则把相关比较写成具体映射。

% 纲要内容（仅作写作计划，尚非教材正文）：
%
% * 锥与柱的显式构造
% * 锥的上同调长正合列
% * 移位和锥的符号检验
%

\subsection{锥与柱的显式构造}
\label{b4:subsec:0602:01}

上同调长正合列从短正合列出发，而许多问题最初只有一个映射 \(f:C\to D\)。要把它转成上同调之间的比较，可以构造一个逐次分裂的复形短列，左端为 \(D\)，右端为 \(C[1]\)。于是中项的次数 \(n\) 取 \(D^n\oplus C[1]^n=D^n\oplus C^{n+1}\)，两个对角微分分别取 \(\mathrm{d}_D^n\)、\(\mathrm{d}_{C[1]}^n=-\mathrm{d}_C^{n+1}\)。再把给定的 \(f^{n+1}:C^{n+1}\to D^{n+1}\) 放在非对角位置，便得到下面的映射锥。移位使这个非对角项的源、目标次数能够相接，所得长正合列将同时比较两端的上同调与 \(H(f)\)。

\begin{definition}\label{b4:def:mapping-cone}
设 \(f:C\to D\) 是加性范畴中的上链复形映射。规定
\[
\operatorname{Cone}(f)^n=D^n\oplus C^{n+1},\qquad
 \mathrm{d}_{\operatorname{Cone}(f)}^n=
 \begin{pmatrix}\mathrm{d}_D^n&f^{n+1}\\0&-\mathrm{d}_C^{n+1}\end{pmatrix}.
\]
所得复形称为 \(f\) 的\term[yingshe-zhui]{映射锥}[mapping cone]。在模中写作 \(\mathrm{d}(y,x)=(\mathrm{d}_Dy+f(x),-\mathrm{d}_Cx)\)。
\end{definition}
\symidx{\(\operatorname{Cone}(f)\)：复形映射的映射锥}

右下角的负号正是\cref{b4:def:complex-shift}中 \(C[1]\) 的微分符号。矩阵平方的非对角项为 \(\mathrm{d}_Df-f\mathrm{d}_C=0\)，对角项为两个微分的平方，因此确为复形。若把右下角的负号删去，非对角项通常会变成 \(\mathrm{d}_Df+f\mathrm{d}_C\)，不能保证为零。

映射柱要把 \(f\) 分解成 \(C\xrightarrow{j}\operatorname{Cyl}(f)\xrightarrow{r}D\)，其中 \(j\) 逐次分裂单，\(r\) 为同伦等价，并保持 \(rj=f\)。为给 \(C\) 留出嵌入，同时保留 \(D\) 的同伦信息，可以在 \(D\) 上添入由 \(C^n\oplus C^{n+1}\) 组成的可缩辅助复形。下面的三个分量正是 \(D\) 与这份辅助数据；定义后将用坐标变换识别它们。

\begin{definition}\label{b4:def:mapping-cylinder}
设 \(f:C\to D\) 是加性范畴中的上链复形映射。规定
\[
\operatorname{Cyl}(f)^n=D^n\oplus C^n\oplus C^{n+1},
\qquad
 \mathrm{d}(y,c,a)=(\mathrm{d}_Dy+f(a),\ \mathrm{d}_Cc-a,\ -\mathrm{d}_Ca).
\]
所得复形称为 \(f\) 的\term[yingshe-zhu]{映射柱}[mapping cylinder]。元组公式在一般加性范畴中表示相应双积矩阵。
\end{definition}
\symidx{\(\operatorname{Cyl}(f)\)：复形映射的映射柱}

这三项可以从一个坐标变换理解。令 \(z=y+f(c)\)，即把 \((y,c,a)\) 换成 \((z,c,a)\)，逆变换为 \((z,c,a)\mapsto(z-f(c),c,a)\)。因为 \(f\) 与微分交换，新坐标中的微分为
\[
\begin{gathered}
\mathrm{d}_Dy+f(a)+f(\mathrm{d}_Cc-a)
=\mathrm{d}_D(y+f(c))=\mathrm{d}_Dz,\\
\mathrm{d}(z,c,a)=(\mathrm{d}_Dz,\,\mathrm{d}_Cc-a,\,-\mathrm{d}_Ca).
\end{gathered}
\]
第一分量因而与后两分量分开：映射柱同构于 \(D\oplus K\)，其中 \(K^n=C^n\oplus C^{n+1}\)，微分为 \((c,a)\mapsto(\mathrm{d}_Cc-a,-\mathrm{d}_Ca)\)。辅助复形 \(K\) 可缩，收缩为 \((c,a)\mapsto(0,-c)\)。映射柱是在 \(D\) 上添入这一份同伦上为零的复形；同时 \(c\mapsto(0,c,0)\) 为 \(C\) 留出逐次分裂的嵌入。在新坐标中，该嵌入是 \(c\mapsto(f(c),c,0)\)，而到 \(D\) 的投影正是原坐标中的 \(y+f(c)\)。

\begin{proposition}\label{b4:prop:cylinder-factorization}
设 \(\mathcal A\) 为加性范畴，\(C,D\) 为其中的上链复形，\(f:C\to D\) 为复形映射。令 \(j:C\to\operatorname{Cyl}(f)\)、\(r:\operatorname{Cyl}(f)\to D\)、\(s:D\to\operatorname{Cyl}(f)\) 的逐次公式为
\[
j(c)=(0,c,0),\qquad r(y,c,a)=y+f(c),\qquad s(y)=(y,0,0).
\]
则 \(rj=f\)、\(rs=1_D\)，\(r,s\) 互为同伦逆。\(j:C\to\operatorname{Cyl}(f)\) 逐次为分裂单态射，且 \(q^n(y,c,a)=(y,a)\) 给出 \(j\) 的余核复形 \(\operatorname{Cone}(f)\)。在交换范畴中因而有逐次分裂短正合列
\[
0\longrightarrow C\xrightarrow{j}\operatorname{Cyl}(f)
\xrightarrow{q}\operatorname{Cone}(f)\longrightarrow0.
\]
\end{proposition}
\begin{proof}
将微分公式再用一次，三个分量分别为 \(\mathrm{d}_D^2y+\mathrm{d}_Df(a)-f\mathrm{d}_C(a)\)、\(\mathrm{d}_C^2c-\mathrm{d}_Ca+\mathrm{d}_Ca\)、\(\mathrm{d}_C^2a\)，均为零。由 \(f\) 与微分交换，\(j,r,s\) 都是复形映射；两个复合等式直接成立。

还需证明 \(1-sr\) 零伦。按同伦约定，比较恒等态射与 \(sr\)，要找到 \(h\) 使 \(1-sr=\mathrm{d}h+h\mathrm{d}\)。定义降次态射 \(h(y,c,a)=(0,0,-c)\)，则
\[
\mathrm{d}h(y,c,a)=(-f(c),c,\mathrm{d}_Cc),\qquad
h\mathrm{d}(y,c,a)=(0,0,-\mathrm{d}_Cc+a),
\]
相加得 \((-f(c),c,a)=(1-sr)(y,c,a)\)。因此 \(sr\simeq1\)。

\(j^n\) 的逐次左逆是 \(\pi_C^n(y,c,a)=c\)，满足 \(\pi_C^nj^n=1_{C^n}\)。逐次投影到第一、第三个双积分量得到 \(q\)，其逐次截面为 \(\sigma^n(y,a)=(y,0,a)\)，满足 \(q^n\sigma^n=1\)。其目标上的微分为 \((\mathrm{d}_Dy+f(a),-\mathrm{d}_Ca)\)，正是映射锥的微分，且 \(qj=0\)。若复形映射 \(u:\operatorname{Cyl}(f)\to E\) 满足 \(uj=0\)，则双积泛性质给出唯一的态射族 \(v^n:\operatorname{Cone}(f)^n\to E^n\)，使 \(u^n=v^nq^n\)。由于每个 \(q^n\) 都是分裂满态射，等式
\[
\begin{aligned}
\bigl(\mathrm{d}_E^nv^n-v^{n+1}\mathrm{d}_{\operatorname{Cone}(f)}^n\bigr)q^n
&=\mathrm{d}_E^nu^n-u^{n+1}\mathrm{d}_{\operatorname{Cyl}(f)}^n\\
&=0
\end{aligned}
\]
可以在右边消去满态射 \(q^n\)，得到 \(\mathrm{d}_E^nv^n=v^{n+1}\mathrm{d}_{\operatorname{Cone}(f)}^n\)。所以逐次分解组成复形映射 \(v\)，\(q\) 满足 \(j\) 在复形范畴中的余核泛性质。
\end{proof}

这里逐次分裂是指每个次数上的短正合列分裂；复形分裂还要求分裂态射与微分交换。上面的 \(\pi_C\) 与 \(\sigma\) 通常不满足这个要求：\(\pi_C\mathrm{d}(y,c,a)=\mathrm{d}_Cc-a\)，而 \(\mathrm{d}_C\pi_C(y,c,a)=\mathrm{d}_Cc\)；\(\mathrm{d}\sigma(y,a)\) 的中间分量为 \(-a\)，而 \(\sigma\mathrm{d}(y,a)\) 的中间分量为零。可缩则是单个复形的恒等映射零伦的条件；这里可缩的是辅助复形 \(K\)，映射柱本身仍保留 \(D\) 的同伦信息。

\subsection{锥的上同调长正合列}
\label{b4:subsec:0602:02}

\begin{proposition}\label{b4:prop:cone-long-exact}
设 \(\mathcal A\) 是交换范畴，\(C,D\) 为其中的上链复形，\(f:C\to D\) 是复形映射。嵌入 \(y\mapsto(y,0)\) 与投影 \((y,x)\mapsto x\) 给出逐次分裂短正合列
\[
0\longrightarrow D\longrightarrow\operatorname{Cone}(f)
\longrightarrow C[1]\longrightarrow0.
\]
其连接态射在 \(H^n(C[1])=H^{n+1}(C)\) 的识别下就是 \(H^{n+1}(f)\)。因此有自然上同调长正合列
\[
\cdots\to H^n(C)\xrightarrow{H^n(f)}H^n(D)
\to H^n\operatorname{Cone}(f)\to H^{n+1}(C)
\xrightarrow{H^{n+1}(f)}H^{n+1}(D)\to\cdots.
\]
\end{proposition}
\begin{proof}
微分矩阵表明嵌入与投影均为复形映射，前者为后者的核，后者为前者的余核；这些泛性质由逐次双积分解得到。先应用\cref{b4:thm:cohomology-long-exact}，所得原始片段为
\[
H^n(D)\longrightarrow H^n\operatorname{Cone}(f)
\longrightarrow H^n(C[1])\xrightarrow{\delta^n}H^{n+1}(D).
\]
再用 \(H^n(C[1])=H^{n+1}(C)\) 识别移位项，并计算 \(\delta^n\)。投影的逐次截面 \(x\mapsto(0,x)\) 不必为复形映射，因为其微分还产生 \(f(x)\) 分量。对余圈子对象 \(Z^{n+1}(C)\) 使用这一截面，再作锥微分，得到进入 \(D^{n+1}\) 的态射 \(f^{n+1}|_{Z^{n+1}(C)}\)。由\cref{b4:prop:connecting-map-construction}，它在上同调上诱导的正是连接态射 \(H^{n+1}(f)\)。在模中，取余圈 \(x\in Z^{n+1}(C)\)，它代表 \(H^n(C[1])\) 中的同一个类。计算为
\[
\begin{aligned}
x&\longmapsto(0,x),\\
\mathrm{d}_{\operatorname{Cone}(f)}^n(0,x)
&=(f^{n+1}(x),-\mathrm{d}_C^{n+1}x)\\
&=(f^{n+1}(x),0),
\end{aligned}
\]
故 \(\delta^n[x]=[f^{n+1}(x)]\)，没有额外负号。最终长列中的前一个箭头 \(H^n(C)\to H^n(D)\) 来自上一次数的连接 \(\delta^{n-1}:H^{n-1}(C[1])\to H^n(D)\)，在移位识别下正是 \(H^n(f)\)。
\end{proof}

\begin{example}\label{b4:exm:cone-single-map}
把左 \(R\)-模 \(M,N\) 放在次数零，同态 \(u:M\to N\) 的锥只有次数 \(-1,0\) 非零，恰为 \(M\xrightarrow{u}N\)。所以
\[
H^{-1}\operatorname{Cone}(u)=\ker u,\qquad
H^0\operatorname{Cone}(u)=\operatorname{coker}u.
\]
例如 \(\mathbb Z\xrightarrow{m}\mathbb Z\) 的锥在 \(m\ne0\) 时只剩零次上同调 \(\mathbb Z/m\mathbb Z\)。因此 \(\operatorname{Cone}(u)\) 零调，当且仅当 \(\ker u=0\)、\(\operatorname{coker}u=0\)，也就是 \(u\) 为模同构。下一节的\cref{b4:prop:cone-quasi-isomorphism}将把这一原型推广到任意复形映射：锥零调恰好检测它是否在所有次数诱导上同调同构。
\end{example}

\subsection{移位和锥的符号检验}
\label{b4:subsec:0602:03}

锥的不同符号约定可以彼此同构，但必须明确写出比较映射。例如有些文献把 \(C^{n+1}\) 放在前面，并在 \(f\) 前加负号；这时交换两个分量并改变其中一个分量的符号，才能与本书定义比较。

\begin{proposition}\label{b4:prop:cone-sign-compatibility}
设 \(\mathcal A\) 为加性范畴，\(C,D\) 为其中的上链复形，\(f,g:C\to D\) 为复形映射。以下元组公式表示相应双积上的态射矩阵。
\begin{enumerate}
\item 若态射族 \(h^n:C^n\to D^{n-1}\) 满足 \(f-g=\mathrm{d}_Dh+h\mathrm{d}_C\)，则次数 \(n\) 上的公式 \((y,x)\mapsto(y+h^{n+1}(x),x)\) 给出 \(\operatorname{Cone}(f)\cong\operatorname{Cone}(g)\)。这个同构由所给的 \(h\) 确定，不同同伦一般会给出不同同构。
\item 对整数 \(r\)，映射 \((y,x)\mapsto(y,(-1)^r x)\) 给出 \(\operatorname{Cone}(f[r])\cong\operatorname{Cone}(f)[r]\)。
\item \(\operatorname{Cone}(1_C)\) 可缩，收缩同伦为 \((y,x)\mapsto(0,y)\)。
\end{enumerate}
\end{proposition}
\begin{proof}
第一项的映射与微分交换等价于
\[
\mathrm{d}_Dy+\mathrm{d}_Dh(x)+g(x)=\mathrm{d}_Dy+f(x)-h(\mathrm{d}_Cx),
\]
这恰为同伦公式；逆映射把 \(h\) 换成 \(-h\)。第二项中 \(f[r]^n=f^{n+r}\)，不另加符号。源 \(\operatorname{Cone}(f[r])^n\) 与目标 \(\operatorname{Cone}(f)[r]^n\) 的底对象都是 \(D^{n+r}\oplus C^{n+r+1}\)，差别只在微分。暂略已经移到 \(n+r\) 的指标，源的微分矩阵为
\[
\begin{pmatrix}(-1)^r\cdot\mathrm{d}_D&f\\0&-(-1)^r\cdot\mathrm{d}_C\end{pmatrix},
\]
目标的微分矩阵为 \((-1)^r\begin{pmatrix}\mathrm{d}_D&f\\0&-\mathrm{d}_C\end{pmatrix}\)。令 \(\epsilon=(-1)^r\)，比较态射 \(t^n=\operatorname{diag}(1,\epsilon)\) 在第二分量补上符号。用 \(\epsilon^2=1\) 计算，两种复合的矩阵均为
\[
t^{n+1}\mathrm{d}_{\operatorname{Cone}(f[r])}^n
=\mathrm{d}_{\operatorname{Cone}(f)[r]}^nt^n
=\begin{pmatrix}\epsilon\cdot\mathrm{d}_D&f\\0&-\mathrm{d}_C\end{pmatrix}.
\]
而 \((t^n)^2=1\)，故给出复形同构。第三项令 \(h(y,x)=(0,y)\)，则 \(\mathrm{d}h(y,x)=(y,-\mathrm{d}y)\)、\(h\mathrm{d}(y,x)=(0,\mathrm{d}y+x)\)，其和为 \((y,x)\)。
\end{proof}

移位后的同伦也须改符号：若 \(h\) 是 \(f\simeq g\) 的同伦，则 \(h_r^n=(-1)^r h^{n+r}\) 是 \(f[r]\simeq g[r]\) 的同伦。这从两边微分均乘 \((-1)^r\) 直接得出。当 \(r\) 为奇数时，若只移动指标而不乘这个符号，同伦公式右边会多出整体负号；\(r\) 为偶数时则无此差异。

% !TeX root = ../../Book4.tex
% 纲要标题：### 7.3 拟同构
% 纲要标签：b4:sec:0603
% 纲要位置：计划纲要.md，第 3258 行。

\section{拟同构}
\label{b4:sec:0603}

同调上的可逆性弱于在复形层面构造同伦逆。映射锥提供拟同构的判据，而非分裂正合列说明这两种可逆性为何会分离。

% 纲要内容（仅作写作计划，尚非教材正文）：
%
% * 拟同构的定义
% * 与同伦等价的区别
% * 非分裂正合列产生的基本例子
%

\subsection{拟同构的定义}
\label{b4:subsec:0603:01}

\begin{definition}\label{b4:def:quasi-isomorphism}
设 \(\mathcal A\) 是交换范畴，\(f:C\to D\) 是复形映射。若每个整数 \(n\) 的 \(H^n(f):H^n(C)\to H^n(D)\) 都是同构，则称 \(f\) 为\term[ni-tonggou]{拟同构}[quasi-isomorphism]。
\end{definition}

拟同构必须是已经给定的复形映射。“两个复形的同调对象逐次同构”只给出对象之间的比较，未必指定了能够实现这些同构的复形映射。更不能由此断言原来的某个映射就是拟同构。

\begin{proposition}\label{b4:prop:cone-quasi-isomorphism}
设 \(\mathcal A\) 是交换范畴，\(f:C\to D\) 是 \(\mathcal A\) 中上链复形之间的映射。则 \(f\) 是拟同构，当且仅当 \(\operatorname{Cone}(f)\) 零调。对 \(\mathcal A\) 中任意可复合的复形映射 \(f:C\to D\)、\(g:D\to E\)，若 \(f,g,gf\) 的任意两个为拟同构，则第三个也是拟同构。
\end{proposition}
\begin{proof}
锥的长正合列中有
\[
H^n(C)\xrightarrow{H^n(f)}H^n(D)\to
H^n\operatorname{Cone}(f)\to H^{n+1}(C)
\xrightarrow{H^{n+1}(f)}H^{n+1}(D).
\]
若左右两个 \(H(f)\) 都是同构，中间对象到右侧的箭头既单又为零，故中间对象为零。若所有锥同调均为零，则相关正合段给出 \(H^n(f)\) 既单又满，故为同构。最后，对 \(H^n(f),H^n(g)\) 使用同构的复合与消去性质，再让 \(n\) 遍历全部整数即可。
\end{proof}

\begin{proposition}\label{b4:prop:short-exact-quasi-comparison}
设交换范畴中两条复形短正合列之间给定态射。如果三个竖直复形映射中有两个为拟同构，则第三个也为拟同构。
\end{proposition}
\begin{proof}
由\cref{b4:thm:cohomology-long-exact}得到两条长正合列及其间交换图。对待证映射 \(H^n\) 所在的项，取其前后各两项组成五项正合图。其余四个竖直态射均是另两个复形在某个次数的同调态射，所以都是同构。\cref{b4:thm:five-lemma}证明中间态射也是同构。
\end{proof}

\subsection{与同伦等价的区别}
\label{b4:subsec:0603:02}

同伦等价必为拟同构，因为同伦逆在同调上成为真正的逆；这是\cref{b4:prop:homotopy-cohomology}的直接应用。拟同构只要求同调上可逆，同伦等价还要求构造反向的复形映射及两个同伦，故要求更强。

\begin{proposition}\label{b4:prop:field-complex-splitting}
设 \(k\) 是域，\(C\) 是任意 \(k\)-向量空间复形。把各个 \(H^n(C)\) 组成微分全零的复形 \(H(C)\)，则 \(C\) 与 \(H(C)\) 同伦等价。因此 \(k\)-向量空间复形之间的每个拟同构都是同伦等价。所用同伦等价一般依赖补空间的选择。
\end{proposition}
\begin{proof}
在每个次数选择补空间
\[
Z^n(C)=B^n(C)\oplus V^n,\qquad C^n=Z^n(C)\oplus W^n.
\]
微分的限制
\[
\mathrm{d}_C^n|_{W^n}:W^n\xrightarrow{\sim}B^{n+1}(C)
\]
是同构，因为其核与 \(W^n\) 的交为零且其像就是全部边界。令 \(q^n:Z^n(C)\to H^n(C)\) 为商映射，则其限制 \(\rho^n=q^n|_{V^n}\) 也是同构。\(H^n(C)\) 本身是商空间；通过 \((\rho^n)^{-1}\)，才能将它识别为选定的代表子空间 \(V^n\subseteq C^n\)。

利用分解 \(C^n=B^n(C)\oplus V^n\oplus W^n\)，定义
\[
i^n:H^n(C)\longrightarrow C^n,\qquad
i^n(\eta)=(\rho^n)^{-1}(\eta),
\]
以及
\[
p^n:C^n\longrightarrow H^n(C),\qquad
p^n(b+v+w)=\rho^n(v).
\]
\(i^n\) 的像在余圈中，而 \(p^{n+1}\) 在微分的像 \(B^{n+1}(C)\) 上为零，故它们组成复形映射 \(i:H(C)\to C\)、\(p:C\to H(C)\)，且 \(pi=1_{H(C)}\)。再定义
\[
h^n(b+v+w)
=\bigl(\mathrm{d}_C^{n-1}|_{W^{n-1}}\bigr)^{-1}(b).
\]
逐个检查这三个分量，得到
\[
\mathrm{d}_C^{n-1}h^n+h^{n+1}\mathrm{d}_C^n
=1_{C^n}-i^np^n.
\]
于是 \(ip\simeq1_C\)，这给出所需同伦等价。

对拟同构 \(f:C\to D\)，给两者分别作上述选择。\(i_C\) 选取同调类的余圈代表，而 \(p_D\) 在余圈上的限制就是取同调类，因此 \(p_Dfi_C=H(f)\) 是微分为零复形之间的同构。取 \(g=i_C H(f)^{-1}p_D\)。由同伦对复合的相容性，
\[
gf\simeq gf i_Cp_C=i_Cp_C\simeq1_C,
\qquad
fg\simeq i_Dp_Dfg=i_Dp_D\simeq1_D.
\]
所以 \(g\) 是 \(f\) 的同伦逆。选择补空间改变了 \(i,p,h\)，所以这不是无须选择的复形分解。
\end{proof}

上述证明使用的是向量空间每个短正合列都分裂。对一般模，核和像可能没有补模；不能只把 \(k\) 改成任意环便保留结论。

\subsection{非分裂正合列产生的基本例子}
\label{b4:subsec:0603:03}

\begin{example}\label{b4:exm:resolution-quasi-not-homotopy}
令 \(P\) 为次数 \(-1,0\) 的复形 \(\mathbb Z\xrightarrow{2}\mathbb Z\)，令 \(M=(\mathbb Z/2)[0]\)。商映射给出增广 \(\varepsilon:P\to M\)，且 \(H^{-1}(P)=0\)、\(H^0(P)=\mathbb Z/2\)，所以它是拟同构。

\ProseParagraphBreak
但任意复形映射 \(g:M\to P\) 的零次分量 \(\mathbb Z/2\to\mathbb Z\) 都为零，故 \(g=0\)。在集中于零次的复形 \(M\) 上，所有可能的降次同伦也为零，因此 \(\varepsilon g=0\) 不可能同伦于 \(1_M\)。所以 \(\varepsilon\) 不是同伦等价。
\end{example}

\begin{example}\label{b4:exm:additive-not-quasi-preserving}
对上一例逐次应用加性函子 \(\operatorname{Hom}_{\mathbb Z}(\mathbb Z/2,-)\)。由于 \(\operatorname{Hom}_{\mathbb Z}(\mathbb Z/2,\mathbb Z)=0\)，所得 \(F(P)\) 是零复形；而 \(F(M)\) 集中于零次且等于 \(\mathbb Z/2\)。因此 \(F(\varepsilon)\) 不是拟同构。
\end{example}

加性函子会保持已经给定的链同伦，因为它保持同伦公式中的加法与复合，却未必保持拟同构。后面先取合适消解再施加函子，正是为了处理这种差别。对非分裂短正合列形成的零调复形，零复形到它的映射也是拟同构；若它不是可缩复形，就同样不会成为同伦等价。

% !TeX root = ../../Book4.tex
% 纲要标题：### 7.4 双复形的预备
% 纲要标签：b4:sec:0604
% 纲要位置：计划纲要.md，第 3264 行。

\section{双复形的预备}
\label{b4:sec:0604}

映射锥把拟同构的判定化为一个复形的正合性。若分别为张量积的两个变量选取消解，两个次数会同时出现；比较这些消解时，知道每列的上同调相同，还须说明合成后的复形是否也有相同的上同调。下面沿每条对角线合并两个次数，并用映射锥把逐列比较归结为有限消元。对角线上有无限多项时，消元可能需要无限步，直和与直积也可能给出不同的上同调。

% 纲要内容（仅作写作计划，尚非教材正文）：
%
% * 横、纵微分的反交换约定
% * 直和总复形与直积总复形
% * 无界情形的差异暂不隐去
% * 在后续谱序列所用反交换约定下取总微分 \(\mathrm{d}=\mathrm{d}_h+\mathrm{d}_v\)；若从交换方格起步，明确在哪个方向加入 \((-1)^p\) 符号后再总化
%
% ---
%

\subsection{横纵微分的反交换约定}
\label{b4:subsec:0604:01}

给两个复形作张量积时，横向应用第一个复形的微分，纵向应用第二个复形的微分。若直接把两个微分相加，平方里还会出现交叉项；反交换条件恰好使这两项抵消。

\begin{definition}\label{b4:def:double-complex}
设 \(\mathcal A\) 是加性范畴。给定 \(\mathcal A\) 中的对象 \(E^{p,q}\)（\(p,q\in\mathbb Z\)）及态射
\[
\mathrm{d}_h:E^{p,q}\to E^{p+1,q},\qquad
 \mathrm{d}_v:E^{p,q}\to E^{p,q+1},
\]
若每个位置都满足 \(\mathrm{d}_h^2=\mathrm{d}_v^2=0\) 且 \(\mathrm{d}_h\mathrm{d}_v+\mathrm{d}_v\mathrm{d}_h=0\)，则称 \(E\) 为\term[shuang-fuxing]{双复形}[double complex]。这里两个方向的方格反交换。对 \(\mathcal A\) 中的双复形 \(E,F\)，若态射族 \(u^{p,q}:E^{p,q}\to F^{p,q}\) 分别与两个方向的微分交换，则称其为双复形映射。
\end{definition}

例如右 \(R\)-模复形 \(C\) 与左 \(R\)-模复形 \(D\) 给出
\[
E^{p,q}=C^p\otimes_R D^q,\qquad
 \mathrm{d}_h=\mathrm{d}_C\otimes1,\qquad \mathrm{d}_v=(-1)^p(1\otimes \mathrm{d}_D).
\]
先横后纵时，纵向符号变为 \((-1)^{p+1}\)；先纵后横时仍为 \((-1)^p\)，所以两次复合之和为零。这个例子会在 Tor 的两个变量比较中使用。

\ProseParagraphBreak
若最初给出的横纵微分 \(\partial_h,\partial_v\) 满足交换关系 \(\partial_h\partial_v=\partial_v\partial_h\)，则取 \(\mathrm{d}_h=\partial_h\)、\(\mathrm{d}_v|_{E^{p,q}}=(-1)^p\partial_v\) 才得到本书意义下的双复形。后文总化时不再另添一次相同的符号。

\subsection{直和总复形与直积总复形}
\label{b4:subsec:0604:02}

\begin{definition}\label{b4:def:total-complex}
设 \(E\) 是交换范畴 \(\mathcal A\) 中的双复形。若相应的可数余积或积存在，对每个 \(n\in\mathbb Z\) 分别规定
\[
\operatorname{Tot}_{\oplus}(E)^n=\bigoplus_{p+q=n}E^{p,q},\qquad
\operatorname{Tot}_{\Pi}(E)^n=\prod_{p+q=n}E^{p,q},
\qquad \mathrm{d}=\mathrm{d}_h+\mathrm{d}_v.
\]
所得复形分别称为 \(E\) 的\term[zhihe-zongfuxing]{直和总复形}[direct-sum total complex]与\term[zhiji-zongfuxing]{直积总复形}[product total complex]。若每条对角线 \(p+q=n\) 仅有有限个非零对象，则二者都是有限双积，典范相同，简记为 \(\operatorname{Tot}(E)\)。
\end{definition}
\symidx{\(\operatorname{Tot}_{\oplus},\operatorname{Tot}_{\Pi}\)：直和与直积总复形}

直和中每个输入分量至多产生两个输出分量；直积中每个输出分量也只依赖两个输入分量，因此公式都定义了态射。反交换关系给出 \(\mathrm{d}^2=\mathrm{d}_h^2+\mathrm{d}_h\mathrm{d}_v+\mathrm{d}_v\mathrm{d}_h+\mathrm{d}_v^2=0\)。一般交换范畴未必有无限直和或直积；当每条对角线有限时，有限双积已经足够。例如支持在 \(p\ge a,q\ge b\) 或支持在 \(p\le a,q\le b\) 的双复形均满足此条件。

\begin{lemma}[有限对角双复形的消元]\label{b4:lem:bounded-double-complex}
设 \(\mathcal A\) 是交换范畴，\(E\) 是其中的双复形，并且对每个 \(n\in\mathbb Z\)，对角线 \(p+q=n\) 仅有有限个非零项。若每个纵列 \((E^{p,\bullet},\mathrm{d}_v)\) 正合，则 \(\operatorname{Tot}(E)\) 正合。若改为每个横行 \((E^{\bullet,q},\mathrm{d}_h)\) 正合，同一结论成立。
\end{lemma}
\begin{proof}
先用模的元组记号写出有限消元过程。固定总次数 \(n\)，设 \(z=(z_p)\) 为余圈，其中 \(z_p\in E^{p,n-p}\)。从最小的可能非零横指标 \(p\) 开始。该位置的总余圈等式为 \(\mathrm{d}_vz_p+\mathrm{d}_hz_{p-1}=0\)，而较小分量已为零，所以 \(\mathrm{d}_vz_p=0\)。纵列正合，故可取 \(y_p\in E^{p,n-p-1}\) 使 \(\mathrm{d}_vy_p=z_p\)。从 \(z\) 减去 \(\mathrm{d}y_p\)，第 \(p\) 个分量消失，至多改变第 \(p+1\) 个分量。

继续沿横指标增加的方向操作。只有总次数 \(n\) 和 \(n-1\) 两条对角线中的位置会被使用，它们的非零位置并集有限，所以过程在有限步后停止。若某一步所需的 \(E^{p,n-p-1}\) 为零，纵列正合已迫使当前余圈分量为零，无须选择。最终 \(z=\mathrm{d}(\sum_p y_p)\)，所以每个余圈是余边。

在一般交换范畴中，从任意满足 \(\mathrm dz=0\) 的态射 \(z:T\to\operatorname{Tot}(E)^n\) 开始。纵列正合，使纵微分诱导的态射 \(E^{p,n-p-1}\to Z^{n-p}(E^{p,\bullet})\) 为满态射，因而可用\cref{b4:lem:generalized-element-calculus}在满态射后提升 \(z_p\) 到 \(y_p\)。每进行一步提升，都把 \(z\) 与先前选出的全部 \(y_p\) 一同预合成到新的测试对象上，并仍用原符号表示。步骤有限，这些满态射的复合仍是一个满态射 \(e:T'\to T\)；最终所得 \(y:T'\to\operatorname{Tot}(E)^{n-1}\) 满足
\[
\mathrm dy=ze.
\]
该引理的正合性判据于是证明总复形正合。横纵互换仍满足反交换关系，对角有限性不变，故得到横行版本。
\end{proof}

\begin{corollary}\label{b4:cor:double-complex-comparison}
设 \(\mathcal A\) 是交换范畴，\(u:E\to F\) 是其中的双复形映射，且 \(E,F\) 的每条对角线 \(p+q=n\) 均只有有限个非零项。若 \(u\) 在每个纵列上是拟同构，则 \(\operatorname{Tot}(u)\) 是拟同构。逐横行拟同构也有同样结论。
\end{corollary}
\begin{proof}
对每个纵列取映射锥，组成双复形
\[
G^{p,q}=F^{p,q}\oplus E^{p,q+1},\quad
 \mathrm{d}_v^G=\begin{pmatrix}\mathrm{d}_v^F&u\\0&-\mathrm{d}_v^E\end{pmatrix},\quad
 \mathrm{d}_h^G=\begin{pmatrix}\mathrm{d}_h^F&0\\0&-\mathrm{d}_h^E\end{pmatrix}.
\]
两个微分的反交换来自原反交换关系及 \(\mathrm{d}_h^Fu=u\mathrm{d}_h^E\)；新双复形仍逐对角有限。每个纵列是拟同构的锥，由\cref{b4:prop:cone-quasi-isomorphism}正合。上一引理使 \(\operatorname{Tot}(G)\) 正合，而其微分矩阵正是 \(\operatorname{Cone}(\operatorname{Tot}(u))\) 的矩阵。再次用锥判据即可。横行版本交换两个指标进行同一构造。
\end{proof}

\subsection{无界总复形的差异}
\label{b4:subsec:0604:03}

每条对角线有限是上一个证明真正使用的条件；总次数固定，并不自动使一次消元只走有限步。下面的例子同时说明两种无限总化为什么不同。

\begin{example}\label{b4:exm:sum-product-total-difference}
设 \(k\) 是域。对每个整数 \(p\ge0\)，在 \((p,-p)\) 放一维空间 \(k e_p\)，在 \((p,1-p)\) 放 \(k f_p\)，其中 \(e_p,f_p\) 为所选基向量，其余位置为零。定义
\[
\mathrm{d}_v(e_p)=f_p,\qquad \mathrm{d}_h(e_p)=-f_{p+1},\qquad
 \mathrm{d}_v(f_p)=\mathrm{d}_h(f_p)=0.
\]
每个非零纵列都是同构 \(k e_p\to k f_p\)，因而正合；两个微分的所有二次复合为零，故满足反交换条件。总复形只有次数零和一。用各一维分量的基向量 \(e_p,f_p\) 表示坐标，两种总微分均写成
\[
(\alpha_0,\alpha_1,\alpha_2,\ldots)\longmapsto
(\alpha_0,\alpha_1-\alpha_0,\alpha_2-\alpha_1,\ldots).
\]
直积情形中，对任意目标坐标序列 \((\beta_p)\)，唯一的原像由 \(\alpha_p=\sum_{i=0}^p\beta_i\) 给出，所以微分是同构，总复形正合。

\ProseParagraphBreak
直和情形中，每个序列须有有限支集。微分仍单，但像恰好是坐标总和为零的有限支集序列：像的坐标和通过相邻相消为零；反过来，若 \(\sum_i\beta_i=0\)，上述部分和 \(\alpha_p\) 最终为零，确实属于直和。坐标求和映射到 \(k\) 是满射，其核正是微分的像，因此
\[
H^0\operatorname{Tot}_{\oplus}(E)=0,\qquad
H^1\operatorname{Tot}_{\oplus}(E)\cong k,
\]
而直积总复形的上同调全零。一个纵列处处正合的双复形，在直和总化后仍可能出现非零上同调。
\end{example}

解差分方程得到的部分和未必有有限支集。例如次数一的基向量 \(f_0\) 在直积中有原像 \((1,1,1,\ldots)\)，在直和中却没有原像。逐列消元所需的无限多个分量因此未必能组成直和总复形中的一个元素。

\subsection{总微分与总化的符号}
\label{b4:subsec:0604:04}

为便于后续计算，把两种起点区别如下。若横纵微分已经反交换，总微分就是 \(\mathrm{d}_h+\mathrm{d}_v\)；若原方格交换，则应先把纵微分改为 \((-1)^p\partial_v\)，总微分才是 \(\partial_h+(-1)^p\partial_v\)。这两句话描述同一个构造，不能把后一个符号重复施加到前一种约定上。

\begin{example}\label{b4:exm:double-square-sign}
设 \(R\) 交换，\(a,b\in R\)。在四个位置 \((0,0),(1,0),(0,1),(1,1)\) 都放 \(R\)，原横箭头均为乘 \(a\)，原纵箭头均为乘 \(b\)。它们交换。把第 \(p\) 列的纵箭头乘 \((-1)^p\)，并把总次数一按 \((1,0),(0,1)\) 排列，总复形成为
\[
R\xrightarrow{t\mapsto(at,bt)}R^2
\xrightarrow{(u,v)\mapsto-bu+av}R.
\]
复合为 \(-bat+abt=0\)。若漏掉负号，复合一般为 \(2abt\)，所以仅知道原方格交换不足以直接相加。该复形是两个两项复形的张量总化，也是后面两元 Koszul 计算中会反复遇到的形式。
\end{example}

当两套消解形成的双复形逐对角有限时，逐列拟同构的映射锥也逐对角有限，有限消元便把每列的正合性传到总复形，从而得到总复形的拟同构。无限对角线上的比较则还取决于总化方式，以及所需的原像能否组成该总复形中的态射。


\begin{chaptersummary}
同伦把两个复形映射的差写成 \(\mathrm{d}h+h\mathrm{d}\)，因此它们在同调上诱导相同映射。可缩意味着恒等映射零伦，比同调消失更强；对零调复形而言，所缺的条件是每个次数的余圈短正合列都分裂。映射锥把核、余核与连接态射的比较合成一条上同调长正合列，拟同构正好对应锥零调。映射柱则把任意映射分成逐次分裂单态射和同伦等价。张量与 Hom 复形在两个变量都保持链同伦，因而保持同伦等价；保持拟同构则不能只靠这些同伦等式。

\ProseParagraphBreak
双复形的总微分依赖横纵反交换的约定。当每条对角线有限时，逐列或逐行正合可通过有限消元传到总复形；同样的论证给出逐列拟同构的总化比较。无界时，提升得到的序列可能没有有限支集，因而直和与直积总化会产生不同同调。这些区别将在消解、导出函子以及后面的谱序列中继续起作用。
\end{chaptersummary}

\BookExercises

\begin{exercise}\label{b4:ex:06-scalar-homotopy}
固定非零整数 \(m\)，令 \(P=(\mathbb Z\xrightarrow{m}\mathbb Z)\) 位于次数 \(-1,0\)。对整数 \(a,b\)，逐次乘 \(a,b\) 得到复形自映射。证明它们同伦当且仅当 \(m\mid(a-b)\)。
\end{exercise}
\begin{solution}
唯一可能非零的同伦分量是 \(h^0:P^0\to P^{-1}\)，设其为乘 \(t\)。两个次数的同伦公式分别给出 \(a-b=tm\) 与 \(a-b=mt\)，是同一条件。故存在同伦恰好等价于整除条件。
\end{solution}

\begin{exercise}\label{b4:ex:06-homology-zero-not-null}
令 \(P=(\mathbb Z\xrightarrow{2}\mathbb Z)\) 位于次数 \(-1,0\)，\(D=\mathbb Z[0][1]\)。定义 \(f:P\to D\) 的 \(-1\) 次分量为恒等，零次分量为零。证明所有 \(H^n(f)\) 为零，但 \(f\) 不零伦。
\end{exercise}
\begin{solution}
按移位约定，\(D=\mathbb Z[0][1]\) 只有次数 \(-1\) 的项 \(\mathbb Z\)，其微分为零。于是复形映射等式成立。源只有 \(H^0=\mathbb Z/2\mathbb Z\)，目标只有 \(H^{-1}=\mathbb Z\)，所以每阶诱导映射都为零。若 \(f=\mathrm{d}h+h\mathrm{d}\)，唯一可能的分量 \(h^0:\mathbb Z\to\mathbb Z\) 为乘 \(t\)，在次数 \(-1\) 的等式要求 \(1=2t\)，不可能。因此同调相同的两个映射未必同伦。
\end{solution}

\begin{exercise}\label{b4:ex:06-tensor-homotopy}
设 \(R\) 为含幺环，\(C,C'\) 是幺性右 \(R\)-模的上链复形，\(D,D'\) 是幺性左 \(R\)-模的上链复形。设 \(f,g:C\to C'\) 由 \(h\) 给出同伦，\(u,v:D\to D'\) 由 \(k\) 给出同伦。在直和总化的张量复形上，写出 \(f\otimes u\) 与 \(g\otimes v\) 之间的同伦。再取 \(R=\mathbb Z\)、\(C=C'=(\mathbb Z\xrightarrow1\mathbb Z)\)、\(D=D'=(\mathbb Z\xrightarrow1\mathbb Z)\)，两个复形均位于次数 \(0,1\)，取 \(f=g=1_C,u=1_D,v=0\)。计算所得收缩，并检验省去第二因子同伦的次数符号会发生什么。
\end{exercise}
\begin{solution}
把差分成
\[
f\otimes u-g\otimes v=(f-g)\otimes u+g\otimes(u-v).
\]
由\cref{b4:prop:tensor-hom-preserve-homotopy}的两个同伦前后复合相应复形映射，得到总次数为 \(-1\) 的同伦
\[
H(x\otimes y)=h(x)\otimes u(y)+(-1)^p g(x)\otimes k(y)
\quad(x\in C^p).
\]
两个同伦等式相加即得 \(\mathrm{d}H+H\mathrm{d}=f\otimes u-g\otimes v\)。

在所给例子中取 \(h=0\)，取 \(k:D^1\to D^0\) 为恒等，其余分量为零。设 \(e_0,e_1\) 与 \(v_0,v_1\) 分别是两复形的标准基。则非零的收缩值为
\[
H(e_0\otimes v_1)=e_0\otimes v_0,\qquad
H(e_1\otimes v_1)=-e_1\otimes v_0.
\]
它满足 \(\mathrm{d}H+H\mathrm{d}=1\)，所以张量复形可缩。若省去 \((-1)^p\)，得到 \(H_0(x\otimes y)=x\otimes k(y)\)，则
\[
(\mathrm{d}H_0+H_0\mathrm{d})(e_0\otimes v_1)
=e_0\otimes v_1+2e_1\otimes v_0,
\]
不等于 \(e_0\otimes v_1\)。第二因子的降次映射越过第一因子时，次数符号保证这些混合项抵消。
\end{solution}

\begin{exercise}\label{b4:ex:06-additive-contractibility}
证明加性函子逐次作用于复形时保持可缩性。对交换群上的函子
\[
F=\operatorname{Hom}_{\mathbb Z}(\mathbb Z/2\mathbb Z,-),
\]
举例说明它未必保持正合复形。
\end{exercise}
\begin{solution}
由 \(1_C=\mathrm{d}h+h\mathrm{d}\)，加性和复合保持性给出 \(1_{F(C)}=F(\mathrm{d})F(h)+F(h)F(\mathrm{d})\)，所以 \(F(h)\) 是收缩同伦。对正合复形 \(\mathbb Z\xrightarrow2\mathbb Z\to\mathbb Z/2\mathbb Z\)（次数 \(0,1,2\)）应用 \(F\)，得到次数 \(0,1,2\) 的复形
\[
0\longrightarrow0\longrightarrow\mathbb Z/2\mathbb Z,
\]
所以 \(H^2(F(C))=\mathbb Z/2\mathbb Z\ne0\)。
\end{solution}

\begin{exercise}\label{b4:ex:06-no-natural-contraction}
设 \(k\) 是域。把 \(0\to k\xrightarrow{x\mapsto(x,0)}k^2\xrightarrow{(x,y)\mapsto y}k\to0\) 放在次数 \(0,1,2\)。证明此复形可缩，但不存在与其所有复形自同构都交换的收缩同伦。
\end{exercise}
\begin{solution}
短列分裂，故可缩。若 \(h\) 是收缩同伦，次数二的等式要求 \(h^2(t)=(at,t)\)，其中 \(a\in k\)。对任意 \(b\in k\)，在中项用剪切 \(u_b(x,y)=(x+by,y)\)，两端用恒等，得到复形自同构。若 \(h\) 与它交换，就需 \(u_bh^2(t)=h^2(t)\)，即 \((a+b)t=at\)。取 \(b=1,t=1\) 矛盾。因此，每个选定的分裂都能给出收缩，却没有一个收缩对所有这些自同构保持不变；收缩存在不意味着可以自然地选取。
\end{solution}

\begin{exercise}\label{b4:ex:06-cone-euler}
设 \(C,D\) 是有限维 \(k\)-向量空间的有界复形，\(f:C\to D\)。证明 \(\chi(\operatorname{Cone}(f))=\chi(D)-\chi(C)\)。若右端为零，能否断言 \(f\) 是拟同构？
\end{exercise}
\begin{solution}
锥的次数 \(n\) 为 \(D^n\oplus C^{n+1}\)，交错求和即得公式。不能反推拟同构：取 \(C=D=k[0]\)、\(f=0\)，两者 Euler 特征相等，但 \(H^0(f)=0:k\to k\) 不可逆，锥在 \(-1,0\) 两次各有一个 \(k\) 同调。
\end{solution}

\begin{exercise}\label{b4:ex:06-cone-residue}
令 \(R=\mathbb Z/12\mathbb Z\)，把乘 \(4\) 看成集中于零次的 \(R\)-模复形之间的映射。计算其锥的同调，并判断该映射是否为拟同构。
\end{exercise}
\begin{solution}
锥为次数 \(-1,0\) 的 \(R\xrightarrow4R\)。核为 \(\{0,3,6,9\}\cong\mathbb Z/4\mathbb Z\)，余核为 \(R/4R\cong\mathbb Z/4\mathbb Z\)，分别是 \(H^{-1}\)、\(H^0\)。锥不正合，故乘四不是拟同构。
\end{solution}

\begin{exercise}\label{b4:ex:06-cone-comparison-integers}
令 \(P=(\mathbb Z\xrightarrow2\mathbb Z)\)、\(Q=(\mathbb Z\xrightarrow6\mathbb Z)\) 均在次数 \(-1,0\)。取 \(f^{-1}=1\)、\(f^0=3\)。写出映射锥并直接计算同调，再与锥长正合列核对。
\end{exercise}
\begin{solution}
\(6\cdot1=3\cdot2\)，故 \(f\) 是复形映射。锥在次数 \(-2,-1,0\) 为
\[
\mathbb Z\xrightarrow{t\mapsto(t,-2t)}\mathbb Z^2
\xrightarrow{(u,v)\mapsto6u+3v}\mathbb Z.
\]
第一映射单，第二映射的核恰为第一映射的像，余核为 \(\mathbb Z/3\mathbb Z\)。另一方面，\(H^0(f):\mathbb Z/2\mathbb Z\to\mathbb Z/6\mathbb Z\) 将一送到三，是单射且余核 \(\mathbb Z/3\mathbb Z\)；锥长正合列给出相同结果。
\end{solution}

\begin{exercise}\label{b4:ex:06-cylinder-homotopies}
设 \(f,g:C\to D\) 为模复形映射。证明 \(f\simeq g\) 当且仅当存在复形映射 \(F:\operatorname{Cyl}(1_C)\to D\)，使 \(F(0,c,0)=f(c)\)、\(F(y,0,0)=g(y)\)。写出 \(F\) 与同伦的关系。
\end{exercise}
\begin{solution}
在 \(\operatorname{Cyl}(1_C)^n=C^n\oplus C^n\oplus C^{n+1}\) 中，第一份 \(C^n\) 是映射 \(1_C\) 的目标，第二份是它的源，第三份属于移位复形。任何这样的分次映射都形如 \(F(y,c,a)=g(y)+f(c)+t(a)\)，其中 \(t\) 降一次。将映射柱的微分代入，与微分交换的额外条件为 \(\mathrm{d}_Dt+t\mathrm{d}_C=g-f\)。因此若 \(h\) 满足 \(f-g=\mathrm{d}_Dh+h\mathrm{d}_C\)，取 \(t=-h\)；反之令 \(h=-t\) 即得同伦。这里第三分量把两个端部映射之间的同伦直接保存下来。
\end{solution}

\begin{exercise}\label{b4:ex:06-tensor-breaks-quasi}
对\cref{b4:exm:resolution-quasi-not-homotopy}中的拟同构 \(P\to(\mathbb Z/2\mathbb Z)[0]\) 张量 \((\mathbb Z/2\mathbb Z)[0]\)。证明所得映射不是拟同构，并解释这与张量保持同伦等价并不冲突。
\end{exercise}
\begin{solution}
\(P\otimes\mathbb Z/2\mathbb Z\) 是次数 \(-1,0\) 的 \(\mathbb Z/2\mathbb Z\xrightarrow0\mathbb Z/2\mathbb Z\)，所以负一次同调为 \(\mathbb Z/2\mathbb Z\)。目标张量积仍集中于零次，负一次同调为零，故不是拟同构。这里 \(\mathbb Z/2\mathbb Z\) 不是平坦的 \(\mathbb Z\)-模：对单射 \(\mathbb Z\xrightarrow2\mathbb Z\) 张量后得到零映射 \(\mathbb Z/2\mathbb Z\to\mathbb Z/2\mathbb Z\)，单射性丢失，产生了额外的负一次同调。原映射没有同伦逆；张量保持给定的同伦等式，并不保证保持任意拟同构。
\end{solution}

\begin{exercise}\label{b4:ex:06-homotopy-not-isomorphism}
给出域 \(k\) 上同伦等价但不逐次同构的两个复形。它们可以有不同的总维数吗？
\end{exercise}
\begin{solution}
取 \(C=(k\xrightarrow1k)\) 和零复形。\(C\) 可缩，所以两者之间的零映射互为同伦逆，但 \(C\) 有两个非零项，无法逐次同构于零复形。总维数分别为二和零。在同伦范畴中，\(C\) 与零对象同构；这种同构只要求复合等于恒等的同伦类，不要求每个次数的对象同构。同伦等价保持同调；对各项均有限维的有界 \(k\)-向量空间复形，它也保持 Euler 特征，但不保持各项维数之和。
\end{solution}

\begin{exercise}\label{b4:ex:06-coprime-square-contraction}
将\cref{b4:exm:double-square-sign}取 \(R=\mathbb Z,a=2,b=3\)。证明总复形可缩，并写出一个收缩同伦。
\end{exercise}
\begin{solution}
复形为 \(\mathbb Z\xrightarrow{(2,3)}\mathbb Z^2\xrightarrow{(-3,2)}\mathbb Z\)，次数 \(0,1,2\)。取 \(h^1(u,v)=-u+v\)、\(h^2(t)=(-t,-t)\)。则 \(h^1\mathrm{d}^0=1\)、\(\mathrm{d}^1h^2=1\)，而
\[
\mathrm{d}^0h^1(u,v)+h^2\mathrm{d}^1(u,v)
=(2(-u+v),3(-u+v))+(3u-2v,3u-2v)=(u,v).
\]
所以 \(\mathrm{d}h+h\mathrm{d}=1\)。互素性允许这里出现整数线性组合等于一。
\end{solution}

\begin{exercise}\label{b4:ex:06-first-quadrant-edge}
设交换范畴中的双复形 \(E\) 支持在 \(p,q\ge0\)，且各纵列只可能在 \(q=0\) 有同调。令 \(Z^{p,0}=\ker(\mathrm{d}_v:E^{p,0}\to E^{p,1})\)。证明这些对象随横微分组成复形，并证明自然嵌入 \(Z^{\bullet,0}\to\operatorname{Tot}(E)\) 是拟同构。
\end{exercise}
\begin{solution}
由反交换，\(\mathrm{d}_v\mathrm{d}_hz=-\mathrm{d}_h\mathrm{d}_vz=0\)，故横微分保持 \(Z^{p,0}\)。把它作为仅有 \(q=0\) 一行的双复形。嵌入在每个纵列上诱导同调同构：零次由 \(E^{p,-1}=0\) 得 \(H^0_v=Z^{p,0}\)，其他次数由假设为零。当总次数 \(n\ge0\) 时，\(p+q=n\) 且 \(p,q\ge0\) 迫使 \(0\le p\le n\)；当 \(n<0\) 时没有非零项。因此第一象限保证逐对角有限，应用\cref{b4:cor:double-complex-comparison}即可。
\end{solution}

\begin{exercise}\label{b4:ex:06-third-quadrant-edge}
设交换范畴中的双复形 \(E\) 支持在 \(p,q\le0\)，且各横行只在 \(p=0\) 可能有同调。证明横同调 \(H_h^0(E^{\bullet,q})\) 随纵微分组成复形，并且从总复形到这一复形的自然投影是拟同构。
\end{exercise}
\begin{solution}
因为 \(E^{1,q}=0\)，零次横同调为 \(E^{0,q}/\operatorname{im}\mathrm{d}_h\)。反交换保证 \(\mathrm{d}_v\) 将横边界送到横边界，所以诱导商上的微分。把这些商置于 \(p=0\) 一列，其他列为零；从 \(E\) 到它的商映射在每个横行上是拟同构。当总次数 \(n\le0\) 时，\(p+q=n\) 且 \(p,q\le0\) 迫使 \(n\le p\le0\)；当 \(n>0\) 时没有非零项。因此第三象限同样逐对角有限，由横行版本的总化比较推论得到结论。
\end{solution}

\begin{exercise}\label{b4:ex:06-infinite-diagonal-size}
设 \(k\) 是域，双复形在每个 \((p,-p)\)、\(p\in\mathbb Z\) 放 \(k\)，其余位置为零，所有微分为零。比较两种总复形及自然映射 \(\operatorname{Tot}_{\oplus}\to\operatorname{Tot}_{\Pi}\)，并指出其为何不是拟同构。
\end{exercise}
\begin{solution}
两者都集中于总次数零，分别为 \(\bigoplus_{p\in\mathbb Z}k\) 与 \(\prod_{p\in\mathbb Z}k\)。自然映射是有限支集序列的包含，不满，例如全一序列不在像中。因此零次同调映射不满，不能成为拟同构。即使没有微分，总化所选的积或余积也会改变对象。
\end{solution}

\begin{exercise}\label{b4:ex:06-hom-homotopies}
固定非零整数 \(m\)，令交换群复形 \(P=(\mathbb Z\xrightarrow{m}\mathbb Z)\) 位于次数 \(-1,0\)，令 \(D=\mathbb Z[0]\)。写出 \(\operatorname{Hom}_{\mathbb Z}^\bullet(P,D)\) 的非零项及微分，计算每个整数 \(r\) 的 \(H^r\operatorname{Hom}_{\mathbb Z}^\bullet(P,D)\)。再用\cref{b4:prop:tensor-hom-preserve-homotopy}识别 \(\operatorname{Hom}_{K(\mathbf{Ab})}(P,D[r])\)：在 \(r=1\) 时用具体复形映射代表同伦类，并说明它何时零伦，核对 Hom 余边与移位同伦之间的符号。
\end{exercise}
\begin{solution}
Hom 复形仅有次数零和一非零，分别来自 \(P^0\to D^0\) 与 \(P^{-1}\to D^0\)。两项均为 \(\mathbb Z\)，而次数零的微分为 \(\mathrm{d}_D\beta-\beta\mathrm{d}_P\)，所以复形为
\[
\mathbb Z\xrightarrow{-m}\mathbb Z
\quad\text{，位于次数 }0,1.
\]
因此零次同调为零，一次同调为 \(\mathbb Z/m\mathbb Z\)，其余次数的同调全零。命题给出的同构说明这些分别也是 \(P\to D[r]\) 的同伦类群。

当 \(r=1\) 时，目标 \(D[1]\) 只有次数 \(-1\) 非零。对整数 \(a\)，令 \(\alpha_a^{-1}:P^{-1}\to D[1]^{-1}\) 为乘 \(a\)，其他分量为零；这给出全部复形映射。两个映射 \(\alpha_a,\alpha_b\) 同伦，当且仅当存在乘 \(t\) 的同伦分量 \(h^0:P^0\to D[1]^{-1}\)，使 \(a-b=tm\)。所以类由 \(a\bmod m\) 决定，且 \(\alpha_a\) 零伦当且仅当 \(m\mid a\)。若 \(a=mt\)，对应的次数零 Hom 元素 \(\beta\) 为乘 \(-t\)，于是 \(\mathrm{d}_{\mathrm{Hom}}\beta=a\)，且 \(h=(-1)^1\beta\) 为乘 \(t\)，与上述零伦同伦一致。
\end{solution}
